\chapter{Child Trauma}
\index{Child Trauma}

\section*{Definition and Overview}
Child trauma refers both to physical injury in children and to the psychological harm that follows frightening or distressing events. Physically, children are at particular risk of injuries from falls, burns, drowning, poisoning, and road accidents. Psychologically, trauma describes the response to events that overwhelm a child's ability to cope, such as serious accidents, violence, abuse, neglect, or sudden loss. Many children show distress after such events and recover with support, but some develop ongoing problems, including post-traumatic stress disorder (PTSD). Both meanings matter: an injured child needs urgent medical care, and a child exposed to a distressing event needs understanding and often professional help.

\section*{Epidemiology}
Injury is the leading cause of death and disability in children over 1 year old in high-income countries. Falls are the most common cause of injury in young children, while road accidents, drowning, and burns cause the most serious harm. Boys are injured more often than girls, and injuries are more common in lower-income households and rural areas. Psychological trauma is also common: a significant minority of children experience a potentially traumatic event, and around 1 in 20 may develop PTSD at some point in childhood. Children exposed to abuse, family violence, or repeated adversity are at greatest risk of long-term difficulties.

\section*{Aetiology and Pathophysiology}
Physical injury in children is closely linked to development: as children gain mobility and curiosity before they understand danger, they encounter falls, burns, poisons, and traffic. Small airways, a proportionally large head, and thinner skulls mean children can be hurt by forces that would harm an adult less. Psychological trauma affects the developing brain's stress systems. When a child faces a threat, the body releases stress hormones; repeated or overwhelming stress can keep these systems on high alert, affecting mood, sleep, memory, and the ability to feel safe. The response depends on the child's age, temperament, and the support available afterwards.

\section*{Clinical Features}
After physical injury, the signs depend on the injury itself:

\begin{itemize}
    \item Pain, swelling, bruising, or an obvious deformity
    \item Cuts or bleeding that is hard to stop
    \item Drowsiness, confusion, vomiting, or a change in behaviour after a head injury
    \item Difficulty breathing after a chest injury
    \item Burns with blistering or charring
\end{itemize}

Signs of psychological distress after a traumatic event include:

\begin{itemize}
    \item Replaying the event in play, dreams, or intrusive thoughts
    \item Avoiding places, people, or reminders of what happened
    \item Being easily startled, jumpy, or on edge
    \item Clinging, nightmares, bedwetting, or regression in younger children
    \item Irritability, anger, or withdrawal in older children
    \item Trouble sleeping or concentrating at school
\end{itemize}

\section*{Differential Diagnosis}
\begin{itemize}
    \item A fracture or sprain compared with a simple bruise in an injured limb
    \item Head injury with concussion compared with a more serious brain injury
    \item Normal post-event upset compared with PTSD, which is diagnosed when distress lasts more than a month and interferes with daily life
    \item Anxiety, depression, or ADHD compared with the effects of trauma
    \item Accidental injury compared with non-accidental injury (child abuse) -- requires a careful and sensitive assessment
\end{itemize}

\section*{When to Seek Medical Care}
Seek urgent care for any significant injury, a suspected fracture, a burn bigger than a small area, a cut that will not stop bleeding, or any head injury with loss of consciousness, vomiting, drowsiness, or confusion. See a doctor if emotional or behavioural symptoms after a distressing event persist for more than a couple of weeks or interfere with school, sleep, or relationships.

\textbf{Contact your local emergency services for:}
\begin{itemize}
    \item Unconsciousness or a child not breathing normally
    \item Severe bleeding or a suspected spinal injury
    \item Difficulty breathing after injury
    \item A fit or seizure after a head injury
    \item Any suspicion that a child is in immediate danger
\end{itemize}

\section*{Diagnosis and Investigations}
\begin{itemize}
    \item \textbf{History:} what happened, when, and how the child has been since
    \item \textbf{Examination:} a full head-to-toe check for bruising, tenderness, and deformity
    \item \textbf{X-rays or CT scan:} to detect fractures or internal injury when suspected
    \item \textbf{Head injury observation:} repeated checks of consciousness, pupils, and balance
    \item \textbf{Psychological assessment:} discussion of the child's fears, sleep, play, and behaviour, often with questionnaires
    \item \textbf{Safety planning:} involvement of child protection services where abuse or neglect is a concern
\end{itemize}

\section*{Management}
\begin{itemize}
    \item \textbf{First aid:} controlling bleeding, cooling burns with running water, and keeping a possibly injured spine still
    \item \textbf{Medical treatment:} splints, plaster, wound care, and surgery where needed
    \item \textbf{Pain relief:} paracetamol or ibuprofen at weight-appropriate doses, or stronger medicines in hospital
    \item \textbf{Head injury observation:} hospital monitoring or clear home-care instructions with red-flag warnings
    \item \textbf{Psychological first aid:} reassurance, a return to routine, and allowing the child to talk or play out what happened
    \item \textbf{Talking therapy:} trauma-focused cognitive behaviour therapy for children with PTSD
    \item \textbf{Family support:} helping parents respond calmly and consistently, and involving school where needed
\end{itemize}

\section*{Complications}
\begin{itemize}
    \item Infection, bleeding, or poor healing of wounds and fractures
    \item Concussion, and rarely more serious brain injury after head trauma
    \item Post-traumatic stress disorder or other anxiety disorders
    \item Depression, behaviour problems, or difficulties at school
    \item Avoidance of activities such as swimming or cycling after an accident
    \item In abused children, long-term effects on development and mental health
\end{itemize}

\section*{Prognosis and Outcomes}
Most childhood injuries heal fully with little or no lasting disability. Concussion symptoms usually resolve within a few weeks with rest and gradual return to activity. Most children who experience a frightening event recover over weeks with support, stability, and reassurance from caring adults. Children who develop PTSD often respond well to trauma-focused therapy. A child's recovery is strongly influenced by the response of parents and caregivers -- calm, supportive care after an event is one of the best predictors of a good outcome.

\section*{Prevention and Self-Care}
\begin{itemize}
    \item Use age-appropriate car seats, seat belts, helmets, and safety gates
    \item Fence pools and supervise children constantly near water
    \item Store medicines, cleaning products, and poisons out of reach and locked away
    \item Fit smoke alarms and use guards around heaters, fires, and hot surfaces
    \item Make safe places for play and supervise playgrounds and roads
    \item After any distressing event, keep routines, be patient, and answer the child's questions honestly
    \item Seek professional help early if distress is not settling
\end{itemize}

\section*{Sources and Further Reading}
The following reputable organizations provide reliable information on child injury prevention, head injury in children, and the effects of trauma on children.
\begin{itemize}
    \item NHS. \textit{Information on accidents, first aid, and injuries in children.} https://www.nhs.uk/conditions/accidents-first-aid/
    \item Centers for Disease Control and Prevention. \textit{Child injury prevention resources.} https://www.cdc.gov/injury/childinjury/index.html
    \item World Health Organization. \textit{Child injury prevention fact sheets.} https://www.who.int/health-topics/injuries
    \item MedlinePlus. \textit{Childhood injuries and safety overview.} https://medlineplus.gov/childsafety.html
    \item MSD Manual. \textit{Injuries in children.} https://www.msdmanuals.com/home/injuries-and-poisoning
\end{itemize}
